% --- Archon physics formalization source begin ---
% archon:physics
% archon:covers PhyXMiniProblems/problem_phyx_mini_0641.lean
% archon:source-report reports/phyx_mini/problem_phyx_mini_0641.source.json

\chapter{Physics problem phyx\_mini\_0641}
\label{ch:PhyXMiniProblems_problem_phyx_mini_0641}

\paragraph{Problem source.}
\#\# Physical scenario

An alpha particle approaches a \$\textasciicircum{}\{197\}Au\$ nucleus with a speed of \$1.50 \textbackslash{}times 10\textasciicircum{}7\$ \$m/s\$. As figure shows, the alpha particle is scattered at a \$49\textasciicircum{}\textbackslash{}circ\$ angle at the slower speed of \$1.49 \textbackslash{}times 10\textasciicircum{}7\$ \$m/s\$.

\#\# Question

With what speed does the \$\textasciicircum{}\{197\}Au\$ nucleus recoil?

\#\# Answer choices

- A: A: \$2.48 \textbackslash{}times 10\textasciicircum{}5\$ \$m/s\$

- B: B: \$2.52 \textbackslash{}times 10\textasciicircum{}5\$ \$m/s\$

- C: C: \$2.56 \textbackslash{}times 10\textasciicircum{}5\$ \$m/s\$

- D: D: \$2.62 \textbackslash{}times 10\textasciicircum{}5\$ \$m/s\$

\#\# Figure caption (auxiliary; use the image as primary evidence)

The image depicts a diagram illustrating the scattering of an alpha particle and a gold nucleus. 

- **Alpha Particle (α):** Represented by a small circle on the left with two plus signs, indicating its charge.
- **Gold Nucleus \textbackslash{}(\textasciicircum{}\{197\}\textbackslash{}text\{Au\}:** Shown as a larger circle with multiple plus signs, representing its positive charge.
- **Trajectories:** A blue curved arrow shows the path of the alpha particle as it approaches and deflects away from the gold nucleus.
- **Angle of Deflection:** The angle between the initial and final paths of the alpha particle is labeled as \textbackslash{}(49\textasciicircum{}\textbackslash{}circ\textbackslash{}).
- **Text:** The labels include "α" for the alpha particle, and "197 Au nucleus" for the gold nucleus.

The diagram conveys a scenario of alpha particle scattering in a physics context, possibly related to experiments like the Rutherford scattering experiment.

\#\# Dataset metadata

Nuclear Physics; Physical Model Grounding Reasoning, Multi-Formula Reasoning

\paragraph{Recorded answer/context.}
B: B: \$2.52 \textbackslash{}times 10\textasciicircum{}5\$ \$m/s\$

\paragraph{Figure/image path.}
/root/proposal\_for\_physic/science-mango/phyx\_mini\_run/phyx\_data/test\_image/641.png

\paragraph{Formalization target.}
create a compiling Lean file with sorry bodies at `PhyXMiniProblems/problem\_phyx\_mini\_0641.lean`. The Lean declarations must preserve the physical quantities, dimensions or dimensional roles, figure labels, governing-law hypotheses, and final relation expressed by this problem.
Use Mathlib/Physlib names found through LeanExplore where available. If a physics API is missing, introduce faithful local abstractions rather than scalar placeholder aliases.

\begin{theorem}[Physics formalization target]
\label{thm:physics:phyx_mini_0641:target}
\lean{PhyXMiniProblems.ProblemPhyXMini0641.problem_phyx_mini_0641}
\uses{def:physics:phyx-mini-0641:phyxminiproblems-problemphyxmini0641-speedinmeterspersecond, def:physics:phyx-mini-0641:phyxminiproblems-problemphyxmini0641-alphagoldscatteringsetup, def:physics:phyx-mini-0641:phyxminiproblems-problemphyxmini0641-matchesscatteringproblemdata, def:physics:phyx-mini-0641:phyxminiproblems-problemphyxmini0641-modelsinitiallystationarygoldtarget, def:physics:phyx-mini-0641:phyxminiproblems-problemphyxmini0641-usesmassnumberapproximation, def:physics:phyx-mini-0641:phyxminiproblems-problemphyxmini0641-matchesprimaryscatteringfigure, def:physics:phyx-mini-0641:phyxminiproblems-problemphyxmini0641-interpretsscatteringfiguregeometry, def:physics:phyx-mini-0641:phyxminiproblems-problemphyxmini0641-satisfiesnonrelativisticmomentumlaws, def:physics:phyx-mini-0641:phyxminiproblems-problemphyxmini0641-isuniqueclosestdisplayedrecoilspeed}
The assigned autoformalize agent should translate this physics problem into Lean declarations in the covered file, with theorem and lemma proof bodies written as `by sorry`.
\end{theorem}
\begin{proof}
This is an autoformalization task, not a proof task. Produce statements that can later be proved by the physics prover without weakening the physical model.
\end{proof}

% --- Archon declaration topology begin ---
\section{Declaration topology}
\begin{definition}[Mass Quantity]
  \label{def:physics:phyx-mini-0641:phyxminiproblems-problemphyxmini0641-massquantity}
  \lean{PhyXMiniProblems.ProblemPhyXMini0641.MassQuantity}
  A nonnegative, unit-independent physical rest mass
\end{definition}
\begin{proof}
  This is the declared model definition of Mass Quantity.
\end{proof}
\begin{definition}[Planar Velocity Quantity]
  \label{def:physics:phyx-mini-0641:phyxminiproblems-problemphyxmini0641-planarvelocityquantity}
  \lean{PhyXMiniProblems.ProblemPhyXMini0641.PlanarVelocityQuantity}
  A unit-independent velocity vector in the plane of the scattering figure
\end{definition}
\begin{proof}
  This is the declared model definition of Planar Velocity Quantity.
\end{proof}
\begin{definition}[Planar Momentum Quantity]
  \label{def:physics:phyx-mini-0641:phyxminiproblems-problemphyxmini0641-planarmomentumquantity}
  \lean{PhyXMiniProblems.ProblemPhyXMini0641.PlanarMomentumQuantity}
  A unit-independent momentum vector in the plane of the scattering figure
\end{definition}
\begin{proof}
  This is the declared model definition of Planar Momentum Quantity.
\end{proof}
\begin{definition}[Diagram Axis]
  \label{def:physics:phyx-mini-0641:phyxminiproblems-problemphyxmini0641-diagramaxis}
  \lean{PhyXMiniProblems.ProblemPhyXMini0641.DiagramAxis}
  The two Cartesian directions used in the supplied diagram
\end{definition}
\begin{proof}
  This is the declared model definition of Diagram Axis.
\end{proof}
\begin{definition}[to Fin]
  \label{def:physics:phyx-mini-0641:phyxminiproblems-problemphyxmini0641-diagramaxis-tofin}
  \lean{PhyXMiniProblems.ProblemPhyXMini0641.DiagramAxis.toFin}
  \uses{def:physics:phyx-mini-0641:phyxminiproblems-problemphyxmini0641-diagramaxis}
  Coordinate index corresponding to a displayed diagram axis
\end{definition}
\begin{proof}
  This is the declared model definition of to Fin.
\end{proof}
\begin{definition}[mass Readout]
  \label{def:physics:phyx-mini-0641:phyxminiproblems-problemphyxmini0641-massreadout}
  \lean{PhyXMiniProblems.ProblemPhyXMini0641.massReadout}
  \uses{def:physics:phyx-mini-0641:phyxminiproblems-problemphyxmini0641-massquantity}
  Read a physical mass in a selected mass unit
\end{definition}
\begin{proof}
  This is the declared model definition of mass Readout.
\end{proof}
\begin{definition}[velocity Component Readout]
  \label{def:physics:phyx-mini-0641:phyxminiproblems-problemphyxmini0641-velocitycomponentreadout}
  \lean{PhyXMiniProblems.ProblemPhyXMini0641.velocityComponentReadout}
  \uses{def:physics:phyx-mini-0641:phyxminiproblems-problemphyxmini0641-planarvelocityquantity, def:physics:phyx-mini-0641:phyxminiproblems-problemphyxmini0641-diagramaxis}
  Read one planar velocity component in selected length and time units
\end{definition}
\begin{proof}
  This is the declared model definition of velocity Component Readout.
\end{proof}
\begin{definition}[speed Readout]
  \label{def:physics:phyx-mini-0641:phyxminiproblems-problemphyxmini0641-speedreadout}
  \lean{PhyXMiniProblems.ProblemPhyXMini0641.speedReadout}
  Read a speed magnitude in selected length and time units
\end{definition}
\begin{proof}
  This is the declared model definition of speed Readout.
\end{proof}
\begin{definition}[momentum Component Readout]
  \label{def:physics:phyx-mini-0641:phyxminiproblems-problemphyxmini0641-momentumcomponentreadout}
  \lean{PhyXMiniProblems.ProblemPhyXMini0641.momentumComponentReadout}
  \uses{def:physics:phyx-mini-0641:phyxminiproblems-problemphyxmini0641-planarmomentumquantity, def:physics:phyx-mini-0641:phyxminiproblems-problemphyxmini0641-diagramaxis}
  Read one planar momentum component in coherent selected base units
\end{definition}
\begin{proof}
  This is the declared model definition of momentum Component Readout.
\end{proof}
\begin{definition}[mass In Kilograms]
  \label{def:physics:phyx-mini-0641:phyxminiproblems-problemphyxmini0641-massinkilograms}
  \lean{PhyXMiniProblems.ProblemPhyXMini0641.massInKilograms}
  \uses{def:physics:phyx-mini-0641:phyxminiproblems-problemphyxmini0641-massquantity, def:physics:phyx-mini-0641:phyxminiproblems-problemphyxmini0641-massreadout}
  Kilogram readout of a physical mass
\end{definition}
\begin{proof}
  This is the declared model definition of mass In Kilograms.
\end{proof}
\begin{definition}[speed In Meters Per Second]
  \label{def:physics:phyx-mini-0641:phyxminiproblems-problemphyxmini0641-speedinmeterspersecond}
  \lean{PhyXMiniProblems.ProblemPhyXMini0641.speedInMetersPerSecond}
  \uses{def:physics:phyx-mini-0641:phyxminiproblems-problemphyxmini0641-speedreadout}
  Metres-per-second readout of a physical speed
\end{definition}
\begin{proof}
  This is the declared model definition of speed In Meters Per Second.
\end{proof}
\begin{definition}[velocity Component In Meters Per Second]
  \label{def:physics:phyx-mini-0641:phyxminiproblems-problemphyxmini0641-velocitycomponentinmeterspersecond}
  \lean{PhyXMiniProblems.ProblemPhyXMini0641.velocityComponentInMetersPerSecond}
  \uses{def:physics:phyx-mini-0641:phyxminiproblems-problemphyxmini0641-planarvelocityquantity, def:physics:phyx-mini-0641:phyxminiproblems-problemphyxmini0641-diagramaxis, def:physics:phyx-mini-0641:phyxminiproblems-problemphyxmini0641-velocitycomponentreadout}
  SI readout of one planar velocity component
\end{definition}
\begin{proof}
  This is the declared model definition of velocity Component In Meters Per Second.
\end{proof}
\begin{definition}[Scattering Body]
  \label{def:physics:phyx-mini-0641:phyxminiproblems-problemphyxmini0641-scatteringbody}
  \lean{PhyXMiniProblems.ProblemPhyXMini0641.ScatteringBody}
  The two physical bodies participating in the encounter
\end{definition}
\begin{proof}
  This is the declared model definition of Scattering Body.
\end{proof}
\begin{definition}[mass Number]
  \label{def:physics:phyx-mini-0641:phyxminiproblems-problemphyxmini0641-scatteringbody-massnumber}
  \lean{PhyXMiniProblems.ProblemPhyXMini0641.ScatteringBody.massNumber}
  \uses{def:physics:phyx-mini-0641:phyxminiproblems-problemphyxmini0641-scatteringbody}
  Nucleon number A of each body in the textbook mass approximation
\end{definition}
\begin{proof}
  This is the declared model definition of mass Number.
\end{proof}
\begin{definition}[atomic Number]
  \label{def:physics:phyx-mini-0641:phyxminiproblems-problemphyxmini0641-scatteringbody-atomicnumber}
  \lean{PhyXMiniProblems.ProblemPhyXMini0641.ScatteringBody.atomicNumber}
  \uses{def:physics:phyx-mini-0641:phyxminiproblems-problemphyxmini0641-scatteringbody}
  Proton number Z, hence the positive nuclear charge number
\end{definition}
\begin{proof}
  This is the declared model definition of atomic Number.
\end{proof}
\begin{definition}[Scattering Phase]
  \label{def:physics:phyx-mini-0641:phyxminiproblems-problemphyxmini0641-scatteringphase}
  \lean{PhyXMiniProblems.ProblemPhyXMini0641.ScatteringPhase}
  Whether a state is before or after the scattering interaction
\end{definition}
\begin{proof}
  This is the declared model definition of Scattering Phase.
\end{proof}
\begin{definition}[Planar Motion State]
  \label{def:physics:phyx-mini-0641:phyxminiproblems-problemphyxmini0641-planarmotionstate}
  \lean{PhyXMiniProblems.ProblemPhyXMini0641.PlanarMotionState}
  \uses{def:physics:phyx-mini-0641:phyxminiproblems-problemphyxmini0641-planarvelocityquantity, def:physics:phyx-mini-0641:phyxminiproblems-problemphyxmini0641-planarmomentumquantity}
  Independent kinematic data of one body in one phase of the encounter
\end{definition}
\begin{proof}
  This is the declared model definition of Planar Motion State.
\end{proof}
\begin{definition}[Planar Direction]
  \label{def:physics:phyx-mini-0641:phyxminiproblems-problemphyxmini0641-planardirection}
  \lean{PhyXMiniProblems.ProblemPhyXMini0641.PlanarDirection}
  Qualitative arrow directions visible in the primary raster
\end{definition}
\begin{proof}
  This is the declared model definition of Planar Direction.
\end{proof}
\begin{definition}[Figure Label]
  \label{def:physics:phyx-mini-0641:phyxminiproblems-problemphyxmini0641-figurelabel}
  \lean{PhyXMiniProblems.ProblemPhyXMini0641.FigureLabel}
  Text and symbolic labels visible in the supplied figure
\end{definition}
\begin{proof}
  This is the declared model definition of Figure Label.
\end{proof}
\begin{definition}[Alpha Gold Scattering Figure]
  \label{def:physics:phyx-mini-0641:phyxminiproblems-problemphyxmini0641-alphagoldscatteringfigure}
  \lean{PhyXMiniProblems.ProblemPhyXMini0641.AlphaGoldScatteringFigure}
  \uses{def:physics:phyx-mini-0641:phyxminiproblems-problemphyxmini0641-scatteringbody, def:physics:phyx-mini-0641:phyxminiproblems-problemphyxmini0641-planardirection, def:physics:phyx-mini-0641:phyxminiproblems-problemphyxmini0641-figurelabel}
  Presentation-level data transcribed from the primary image
\end{definition}
\begin{proof}
  This is the declared model definition of Alpha Gold Scattering Figure.
\end{proof}
\begin{definition}[Alpha Gold Scattering Setup]
  \label{def:physics:phyx-mini-0641:phyxminiproblems-problemphyxmini0641-alphagoldscatteringsetup}
  \lean{PhyXMiniProblems.ProblemPhyXMini0641.AlphaGoldScatteringSetup}
  \uses{def:physics:phyx-mini-0641:phyxminiproblems-problemphyxmini0641-massquantity, def:physics:phyx-mini-0641:phyxminiproblems-problemphyxmini0641-scatteringbody, def:physics:phyx-mini-0641:phyxminiproblems-problemphyxmini0641-scatteringphase, def:physics:phyx-mini-0641:phyxminiproblems-problemphyxmini0641-planarmotionstate, def:physics:phyx-mini-0641:phyxminiproblems-problemphyxmini0641-alphagoldscatteringfigure}
  All independent quantities in the scattering model. In particular, the gold recoil speed is an unconstrained dimensionful field here; it is neither defined from the recorded answer nor fixed by any answer-choice value
\end{definition}
\begin{proof}
  This is the declared model definition of Alpha Gold Scattering Setup.
\end{proof}
\begin{definition}[Matches Scattering Problem Data]
  \label{def:physics:phyx-mini-0641:phyxminiproblems-problemphyxmini0641-matchesscatteringproblemdata}
  \lean{PhyXMiniProblems.ProblemPhyXMini0641.MatchesScatteringProblemData}
  \uses{def:physics:phyx-mini-0641:phyxminiproblems-problemphyxmini0641-speedinmeterspersecond, def:physics:phyx-mini-0641:phyxminiproblems-problemphyxmini0641-alphagoldscatteringsetup}
  The two numerical alpha-particle speeds stated in the prose
\end{definition}
\begin{proof}
  This is the declared model definition of Matches Scattering Problem Data.
\end{proof}
\begin{definition}[Models Initially Stationary Gold Target]
  \label{def:physics:phyx-mini-0641:phyxminiproblems-problemphyxmini0641-modelsinitiallystationarygoldtarget}
  \lean{PhyXMiniProblems.ProblemPhyXMini0641.ModelsInitiallyStationaryGoldTarget}
  \uses{def:physics:phyx-mini-0641:phyxminiproblems-problemphyxmini0641-velocitycomponentinmeterspersecond, def:physics:phyx-mini-0641:phyxminiproblems-problemphyxmini0641-alphagoldscatteringsetup}
  The source does not literally say that the gold nucleus is initially at rest, but that standard target-at-rest idealization is necessary to determine the recoil from the supplied data. It is therefore a separate, visible modeling assumption rather than being presented as a source or figure readout
\end{definition}
\begin{proof}
  This is the declared model definition of Models Initially Stationary Gold Target.
\end{proof}
\begin{definition}[Uses Mass Number Approximation]
  \label{def:physics:phyx-mini-0641:phyxminiproblems-problemphyxmini0641-usesmassnumberapproximation}
  \lean{PhyXMiniProblems.ProblemPhyXMini0641.UsesMassNumberApproximation}
  \uses{def:physics:phyx-mini-0641:phyxminiproblems-problemphyxmini0641-massreadout, def:physics:phyx-mini-0641:phyxminiproblems-problemphyxmini0641-massinkilograms, def:physics:phyx-mini-0641:phyxminiproblems-problemphyxmini0641-scatteringbody, def:physics:phyx-mini-0641:phyxminiproblems-problemphyxmini0641-alphagoldscatteringsetup}
  The standard mass-number approximation used by this textbook calculation: each nuclear mass is its nucleon number times one common atomic mass unit. This fixes only the mass ratio 4 : 197, not the requested recoil speed
\end{definition}
\begin{proof}
  This is the declared model definition of Uses Mass Number Approximation.
\end{proof}
\begin{definition}[Matches Primary Scattering Figure]
  \label{def:physics:phyx-mini-0641:phyxminiproblems-problemphyxmini0641-matchesprimaryscatteringfigure}
  \lean{PhyXMiniProblems.ProblemPhyXMini0641.MatchesPrimaryScatteringFigure}
  \uses{def:physics:phyx-mini-0641:phyxminiproblems-problemphyxmini0641-alphagoldscatteringsetup}
  Labels, stylized positive-charge marks, paths, qualitative directions, and the printed angle read from image 641. This structure records no velocity components and no recoil magnitude
\end{definition}
\begin{proof}
  This is the declared model definition of Matches Primary Scattering Figure.
\end{proof}
\begin{definition}[Interprets Scattering Figure Geometry]
  \label{def:physics:phyx-mini-0641:phyxminiproblems-problemphyxmini0641-interpretsscatteringfiguregeometry}
  \lean{PhyXMiniProblems.ProblemPhyXMini0641.InterpretsScatteringFigureGeometry}
  \uses{def:physics:phyx-mini-0641:phyxminiproblems-problemphyxmini0641-speedinmeterspersecond, def:physics:phyx-mini-0641:phyxminiproblems-problemphyxmini0641-velocitycomponentinmeterspersecond, def:physics:phyx-mini-0641:phyxminiproblems-problemphyxmini0641-alphagoldscatteringsetup}
  Cartesian interpretation of the primary figure. The green recoil arrow fixes only the signs of its velocity components; it does not fix their magnitude
\end{definition}
\begin{proof}
  This is the declared model definition of Interprets Scattering Figure Geometry.
\end{proof}
\begin{definition}[Satisfies Nonrelativistic Momentum Laws]
  \label{def:physics:phyx-mini-0641:phyxminiproblems-problemphyxmini0641-satisfiesnonrelativisticmomentumlaws}
  \lean{PhyXMiniProblems.ProblemPhyXMini0641.SatisfiesNonrelativisticMomentumLaws}
  \uses{def:physics:phyx-mini-0641:phyxminiproblems-problemphyxmini0641-diagramaxis, def:physics:phyx-mini-0641:phyxminiproblems-problemphyxmini0641-massreadout, def:physics:phyx-mini-0641:phyxminiproblems-problemphyxmini0641-velocitycomponentreadout, def:physics:phyx-mini-0641:phyxminiproblems-problemphyxmini0641-speedreadout, def:physics:phyx-mini-0641:phyxminiproblems-problemphyxmini0641-momentumcomponentreadout, def:physics:phyx-mini-0641:phyxminiproblems-problemphyxmini0641-scatteringbody, def:physics:phyx-mini-0641:phyxminiproblems-problemphyxmini0641-scatteringphase, def:physics:phyx-mini-0641:phyxminiproblems-problemphyxmini0641-alphagoldscatteringsetup}
  The nonrelativistic laws used by the solution, stated in arbitrary coherent units: speed is the Euclidean magnitude of velocity, p = m v componentwise, and total planar linear momentum is conserved. These are general governing relations and contain no numerical recoil-speed conclusion
\end{definition}
\begin{proof}
  This is the declared model definition of Satisfies Nonrelativistic Momentum Laws.
\end{proof}
\begin{definition}[Answer Choice]
  \label{def:physics:phyx-mini-0641:phyxminiproblems-problemphyxmini0641-answerchoice}
  \lean{PhyXMiniProblems.ProblemPhyXMini0641.AnswerChoice}
  Labels of the four answer choices printed with the problem
\end{definition}
\begin{proof}
  This is the declared model definition of Answer Choice.
\end{proof}
\begin{definition}[speed In Meters Per Second]
  \label{def:physics:phyx-mini-0641:phyxminiproblems-problemphyxmini0641-answerchoice-speedinmeterspersecond}
  \lean{PhyXMiniProblems.ProblemPhyXMini0641.AnswerChoice.speedInMetersPerSecond}
  \uses{def:physics:phyx-mini-0641:phyxminiproblems-problemphyxmini0641-speedinmeterspersecond, def:physics:phyx-mini-0641:phyxminiproblems-problemphyxmini0641-answerchoice}
  Recoil speed printed beside each answer, in metres per second
\end{definition}
\begin{proof}
  This is the declared model definition of speed In Meters Per Second.
\end{proof}
\begin{definition}[recorded Dataset Answer]
  \label{def:physics:phyx-mini-0641:phyxminiproblems-problemphyxmini0641-recordeddatasetanswer}
  \lean{PhyXMiniProblems.ProblemPhyXMini0641.recordedDatasetAnswer}
  \uses{def:physics:phyx-mini-0641:phyxminiproblems-problemphyxmini0641-answerchoice}
  The answer label recorded by the source dataset
\end{definition}
\begin{proof}
  This is the declared model definition of recorded Dataset Answer.
\end{proof}
\begin{definition}[Is Unique Closest Displayed Recoil Speed]
  \label{def:physics:phyx-mini-0641:phyxminiproblems-problemphyxmini0641-isuniqueclosestdisplayedrecoilspeed}
  \lean{PhyXMiniProblems.ProblemPhyXMini0641.IsUniqueClosestDisplayedRecoilSpeed}
  \uses{def:physics:phyx-mini-0641:phyxminiproblems-problemphyxmini0641-speedinmeterspersecond, def:physics:phyx-mini-0641:phyxminiproblems-problemphyxmini0641-alphagoldscatteringsetup, def:physics:phyx-mini-0641:phyxminiproblems-problemphyxmini0641-answerchoice}
  A displayed choice is uniquely closest to the modeled gold recoil speed
\end{definition}
\begin{proof}
  This is the declared model definition of Is Unique Closest Displayed Recoil Speed.
\end{proof}
% --- Archon declaration topology end ---
% --- Archon physics formalization source end ---
